% Transcription par lecture visuelle directe, sans OCR.
% Statut : transcrit ; relecture indépendante à effectuer.
% Source : PDF original, vue 141, côté gauche, page manuscrite 282, après le chapitre douzième.
\chapter{Quatrième suite et fin de la préfecture de M. Franchet. Restauration de la Congrégation. Préfecture de M. Bertrand du Coin.}

L'autorité du Roi ayant été solennellement reconnue et
% Source : PDF original, vue 141, côté droit, page manuscrite 283.
proclamée à Lyon le Samedi saint 9 avril, le gouvernement de la Congrégation ne vit dès lors aucun obstacle à exécuter à l'avenir toutes les dispositions du règlement avec une scrupuleuse exactitude.

\section*{Première assemblée générale}
% Mention marginale de la page manuscrite 283.
Aussi pour commencer à y conformer son administration, le vice Préfet s'empressa de convoquer une assemblée générale pour le premier dimanche de Mai. Il engagea à s'y rendre, non seulement les membres de la Congrégation qui lui étaient constamment demeurés fidèles, mais encore ceux d'entre eux qui par frayeur ou par suite d'une position spéciale avait cru devoir s'en écarter pour un temps. Ces amis saisirent avec empressement cette occasion de se rattacher à une société dont leurs cœurs n'avaient jamais cessé de faire partie. Quant à ceux qui n'avaient jamais abandonné les drapeaux, ils étaient si heureux de pouvoir enfin se réunir sans entraves et sans inquiétude, qu'ils ne se firent pas prier pour accourir au rendez-vous. L'assemblée fut donc nombreuse, du moins eu égard à l'état où se trouvait alors la Congrégation. Si mes souvenirs sont fidèles, 35 ou 36 membres étaient présents.

La section du zèle avait préparé une réception ; ainsi pour la première fois depuis trois ans, un récipiendaire allait être présenté en assemblée générale.

Nous eussions bien désiré que nos confesseurs de la foi eussent pu être de retour pour assister à cette assemblée ; nous n'avions pas encore de leurs nouvelles au moment où l'assemblée fut convoquée, mais leur mise en liberté ne pouvait pas faire l'ombre d'un doute ; nous reçûmes heureusement la nouvelle positive de leur délivrance assez à temps pour pouvoir donner aux membres réunis un nouveau motif de joie en proclamant cette
% Source : PDF original, vue 142, côté gauche, page manuscrite 284.
grande nouvelle dans l'assemblée générale.

Le Vice Préfet remplaça dans cette assemblée soit la lecture du procès verbal, soit l'allocution qu'il est d'usage que le Président adresse au récipiendaire, par un discours dans lequel il déroulait aux yeux de la Congrégation toute la suite de son histoire. Cela lui donna occasion de signaler l'union intime par laquelle il a plu à la divine Providence d'attacher notre société à la Sainte Église de J. C. et à son chef, en lui faisant partager avec cette Église, comme vous êtes à portée de l'observer vous-même, Messieurs, en lisant cet écrit, et les biens et les maux. Formée au moment où la religion se rétablit en France, réduite à se dérober à tous les regards et à pleurer et gémir dans le secret, lorsque l'Église est opprimée par la tyrannie, persécutée dans la personne de ses membres avec le Souverain Pontife, c'est encore avec l'Église et avec son chef que la Congrégation, pénétrée de reconnaissance entonne l'hymne de triomphe lorsqu'il plaît au Seigneur de déployer la force de son bras et de nous rendre enfin la paix en la rendant à l'Église. Que de motifs pour la Congrégation de se montrer plus que jamais inébranlable dans la foi et l'unité Catholique, et en usant de la sécurité de réchauffer son zèle à la vue de tant de merveilles et de se préparer à travailler à la gloire de Dieu avec une nouvelle ardeur.

C'est en plaçant sous les yeux de ce nouveau soldat qui venait s'enrôler sous les drapeaux de la mère de Dieu, le tableau admirable de tout ce que notre bon maître venait d'opérer en faveur de l'Église, de la France et de la Congrégation, que le vice Préfet s'efforçait de dilater en lui ces doux sentiments d'amour et de
% Source : PDF original, vue 142, côté droit, page manuscrite 285.
reconnaissance envers Dieu et envers Marie notre Sainte protectrice ; il parlait bien de l'abondance du cœur, car comme tous les frères il surabondait de joie, et avait peine à contenir l'émotion si vive et si douce qu'excitaient en nous tous, les événements prodigieux dont nous étions alors les témoins.

Nous aurions été appelés à déterminer nous-même l'époque heureuse de cette résurrection qui venait rendre à la vie notre Congrégation qu'il nous eut été impossible de mieux la choisir. Cette bienheureuse assemblée qui versait dans nos cœurs des consolations si bien senties, avait lieu précisément le troisième dimanche après Pâques ; c'était donc le dimanche suivant que revenait le jour consacré par la Congrégation à de solennelles actions de grâce. Certes, depuis que la Congrégation la célébrait, jamais cette fête n'avait pu se présenter dans un moment plus opportun. Le Vice Préfet mit le comble à l'allégresse générale, fit tressaillir tous les cœurs, en annonçant que toutes les mesures étaient prises pour que cette solennité fût accompagnée de tout l'appareil avec lequel elle avait été célébrée dans les plus beaux jours de la Congrégation, et que déjà nous nous étions assurés de toutes les permissions nécessaires pour pouvoir nous remettre en ce jour en possession de cette ancienne chapelle des Chartreux de laquelle nous avons conservé de si doux souvenirs.

Je crois qu'il est inutile d'ajouter ici qu'on se rendit avec un vif empressement à cette réunion pieuse. Vous savez que rien ne fait mieux comprendre la valeur des biens qu'on possédait jadis, comme la privation qu'on éprouve quand on les a perdus. Aussi la Congrégation avait-elle
% Source : PDF original, vue 143, côté gauche, page manuscrite 286.
une véritable soif de retrouver ces jouissances des premières années ; les anciens n'en parlaient qu'avec émotion. Ceux qui jusqu'ici n'avaient connu que par ouï dire le bonheur ineffable dont ces jours de fête remplissaient nos âmes, soupiraient depuis longtemps après l'heureux instant auquel ils pourraient à leur tour en savourer les douceurs ; d'ailleurs nos cœurs avaient été tellement ouverts à la reconnaissance pour toutes les faveurs répandues sur nous avec une immense largesse, qu'ils éprouvaient le besoin de donner aux sentiments dont ils étaient pleins une manifestation éclatante ; aussi ce fut avec enthousiasme que fut entonné le \textit{Te Deum}\footnote{(Note du transcripteur) Littéralement, d'après l'ouverture de l'hymne : « Nous vous louons, ô Dieu ».}, cette prière d'actions de grâce qui réunissait toutes nos pensées et toutes nos affections, s'élevait vers le ciel comme un encens d'agréable odeur. Elle allait remercier pour nous la divine providence pour la paix qui nous était enfin complètement rendue, paix à l'Église, paix à la France, paix à la Congrégation ; triomphe du Souverain Pontife, retour de notre roi, délivrance de notre Préfet et de cet autre ami qui comme lui appuyé sur la croix de J. C. avait été avec lui jugé digne de marcher à la suite du Souverain maître dans la voie des tribulations.

\section*{Fête d'actions de grâce}
% Mention marginale de la page manuscrite 286 ; placée après le récit du Te Deum.
À ces chants d'allégresse et de reconnaissance succéda l'\textit{Exaudiat}, ce psaume de David spécialement consacré par l'Église à prier pour le chef de l'État. Il fut chanté avec tout l'entrain de Français fidèles à leur Roi et heureux de l'avoir retrouvé.

C'est peut-être ici le lieu de placer une remarque importante. La Congrégation était trop intéressée de toute manière à la
% Source : PDF original, vue 143, côté droit, page manuscrite 287.
restauration pour ne pas l'accueillir avec enthousiasme. Elle savait d'ailleurs qu'un enfant fidèle à l'Église est toujours en même temps fidèle à son Roi légitime. Ses membres furent donc unanimes dans leurs sentiments d'amour pour leur roi, et la Congrégation elle-même mit un vif empressement à lui apporter comme un hommage tout filial le tribut habituel public ou particulier de ses prières. Mais nous devons le déclarer formellement pour répondre à ceux qui nous ont accusés et qui nous accusent peut-être encore d'avoir de vue politique ; là s'est constamment borné toute la part que la Congrégation a prise aux affaires de l'État.

Formée par suite de considérations bien plus élevées que peuvent jamais l'être les pauvres et mesquines affaires de ce monde, si constamment soumise aux puissances elle a toujours rendu à César ce qui appartenait à César, elle n'a surtout pas oublié que la gloire de Dieu, l'honneur de Marie et la sanctification réciproque de ses membres étaient l'unique but que s'étaient proposé les fondateurs, et que par conséquent, c'est à rendre à Dieu ce qui est à Dieu que devaient tendre tous ses efforts. C'est pour cela seulement et pour aucun autre motif qu'elle a été instituée. Elle se considère comme trop honorée par cette noble destination pour jamais s'en écarter.

Monsieur, frère du Roi, ayant pris les rênes du gouvernement en attendant l'arrivée de Louis 18, envoya dans les provinces des sujets dévoués, munis de pouvoirs très étendus et avec le titre de commissaires extraordinaires du Roi pour y établir l'autorité royale. La Providence permit que cette
% Source : PDF original, vue 144, côté gauche, page manuscrite 288.
haute mission fût confiée pour Lyon, à notre cher et fidèle associé M. le C\up{te} Alexis de Noailles. Il eut lui-même l'heureuse pensée de choisir notre digne Préfet, M. Franchet, à peine sorti de Ste Pélagie, pour le seconder dans cette mission importante et il l'amena donc en qualité de secrétaire avec lui à Lyon. Ils y arrivèrent inopinément. Le vice Préfet put enfin rendre compte de sa gestion et remettre la Congrégation entre les mains de son chef, si glorieusement couvert de l'auréole de la persécution.

\section*{Arrivée de MM. Franchet, du Coin et de Noailles}
% Mention marginale de la page manuscrite 288.
Tous les membres de la Congrégation éprouvaient un désir bien vif et bien naturel de revoir leur Préfet. Afin d'avancer pour tous ce moment si heureux, on se hâta de convoquer une assemblée générale extraordinaire. M. du Coin aussi de retour dans la patrie depuis quelques jours, s'empressa de s'y rendre et pour que la fête fût complète M. de Noailles voulut bien laisser de côté pendant quelques instants les hautes fonctions dont il était revêtu et l'entourage brillant dont elles étaient environnées, pour se rappeler seulement son titre de Congréganiste associé qu'il savait si bien apprécier et venir en cette qualité se mêler à ses frères qui de leur côté étaient bien heureux de le revoir. Les paroles me manquent ici pour vous faire comprendre tout ce que la tendresse fraternelle faisait éprouver à nos cœurs. Avec quels doux épanchements on s'embrassait, on se félicitait et surtout on remerciait ensemble le Dieu de toute bonté et notre bonne mère, dont plus que jamais nous pouvions apprécier l'assistance tutélaire ! Nous ne nous lassions pas de contempler au milieu de nous ces trois frères qui après
% Source : PDF original, vue 144, côté droit, page manuscrite 289.
une si longue et si douloureuse absence nous étaient enfin rendus. Tous trois avaient souffert pour la cause de la religion. Tous trois avaient mérité le titre de confesseurs de la foi, M. Franchet et du Coin, par la prison ; M. de Noailles par l'exil. En terminant cette délicieuse réunion nous donnâmes à tous les trois le baiser de charité ; nous étions presque hors de nous-même en les serrant tour à tour dans nos bras, au milieu de l'effusion des plus doux sentiments.

\section*{3\up{e} dimanche de juillet 1814}
% Mention marginale de la page manuscrite 289 ; placée avant l'annonce de la fête.
Une autre fête approchait. Oh ! comme elle faisait d'avance palpiter tous nos cœurs ! Vous avez déjà deviné qu'il s'agit ici du troisième dimanche de juillet. Que de choses nous avions à dire en ce jour au cœur de notre mère dont nous allions célébrer la fête. Elle nous avait si bien montré toute sa tendresse maternelle ! Nous pourrions de nouveau donner un libre cours à nos actions de grâces. Nous pouvions rendre à cette solennité tout l'éclat dont elle avait autrefois brillé. Les cœurs étaient tout préparés à s'abandonner à la joie la plus douce. Aussi vous devez penser dans l'état d'exaltation où tous les miracles dont nous étions les témoins mettaient tous les amis de la religion, et par conséquent tous les Congréganistes, avec quel pieux enthousiasme cette solennité fut célébrée.

Dès le matin une grand'messe fut chantée dans notre chapelle des Chartreux ; elle fut accompagnée du \textit{Te Deum} et de tous les autres accessoires qui étaient déjà consacrés par nos usages. Le soir la réunion eut lieu dans une maison de campagne. Là les Congréganistes se livrèrent tout entiers, avec un bien doux laisser aller à tous les sentiments religieux
% Source : PDF original, vue 145, côté gauche, page manuscrite 290.
nationaux et fraternels qui remplissaient leurs âmes. En relisant dans le registre des procès verbaux ceux de l'époque et surtout celui de cette fête, vous pourrez vous faire quelque idée de l'enthousiasme qui dominait alors toutes nos pensées, et vous ne serez plus étonnés qu'aux acclamations de Vive Jésus, vive Marie, vinssent aussi se mêler celles de Vive le Roi, et qu'à notre ancien cri de guerre, De tous les temps, allons toujours en avant, on en joignît un autre qui s'échappait aussi de tous les cœurs, Dieu et le Roi ! Ce mélange quoiqu'il semblât allier en quelque sorte le sacré avec le profane, était cependant dans l'ordre et pouvait très bien se justifier ; un cœur vraiment chrétien accomplit avec zèle tous ses devoirs ; car ils ont tous Dieu pour principe et pour fin, et l'amour de la patrie peut être compté au nombre des vertus.

\section*{Départ de M. Franchet}
% Mention marginale de la page manuscrite 290.
Le moment allait arriver où la Congrégation aurait à élire son Préfet. La mission de M. de Noailles était terminée, il avait quitté notre ville ; M. Franchet s'étant décidé à entrer dans la carrière politique où plus tard il devait occuper un rang distingué, s'était également vu obligé de se séparer de nous. Dieu nous rappelle sans cesse que nous n'avons pas ici bas une demeure permanente ; si quelquefois il veut bien nous accorder quelques moments de jouissance dans ce monde, il a soin qu'elles ne soient que passagères, afin que nos cœurs n'ayent pas le temps de s'y attacher et n'y rencontrent pas le danger de perdre la pensée continuelle du bonheur de l'éternité qui seule est faite pour remplir notre âme ; aussi il nous avait fait retrouver notre ami, il nous l'avait montré pour un temps, et maintenant il
% Source : PDF original, vue 145, côté droit, page manuscrite 291.
nous est enlevé pour toujours, sans que la Congrégation puisse conserver la moindre espérance de le revoir jamais à sa tête. Du moins en la quittant de corps, ce vaillant ami y demeurait plus que jamais attaché de cœur ; jamais, dans quelque position où il se soit trouvé placé, il n'a oublié une société qui lui est si chère, et s'il n'y a pas pu prendre une part active à ses travaux, il a continué à l'édifier de loin par ses bons exemples.

Vous vous rappelez cette société de bonnes œuvres dont nous avons déjà parlé, que M. de Noailles avait fondé à Paris. M. Franchet en devint bientôt un membre distingué. Il fut appelé à y présider la section des Savoyards, fonction dont il s'est honorablement acquitté jusqu'au moment où ses éminentes qualités ont été appréciées à tel point qu'elles ont déterminé le Gouvernement à lui confier la direction de la police générale. C'est bien là qu'en exécutant à la lettre le vœu que forme notre règlement il a répandu autour de lui, disons mieux, les relations étendues que les devoirs de sa place l'avaient forcé d'établir, ont été cause qu'il a répandu au loin la bonne odeur de J. C., par la conduite éminemment noble et désintéressée dont il a donné l'exemple.

Plus la vertu cherchait à se cacher sous le voile de la modestie, plus elle attirait sur lui l'estime générale et cela même de la part de ceux que leur ligne politique avait placés dans le parti opposé.

Nous pouvons l'affirmer, car nous-même en avons été témoins. Loin de décroître en mérite en arrivant aux honneurs, comme hélas ! cela n'arrive que trop souvent à la foiblesse humaine, M. Franchet bien au contraire n'a fait qu'avancer à grands
% Source : PDF original, vue 146, côté gauche, page manuscrite 292.
pas dans la carrière de la vertu. C'est que M. Franchet est essentiellement un homme intérieur ; il était bien pénétré de cette vérité, qu'il n'était qu'un instrument entre les mains du Souverain maître, et jamais il ne s'est proposé autre chose que sa gloire.

À ces vertus chrétiennes il joignait encore cette capacité, ce coup d'œil juste et ce sang froid qui constituent le véritable homme d'état.

Quel bonheur pour la Congrégation d'avoir pu compter dans ses rangs des hommes tels que MM. Franchet et du Coin. En plaçant sous nos yeux de semblables modèles le Seigneur dans sa bonté a voulu ajouter encore pour nous à tous les autres moyens de salut dont il nous a entourés, ceux que nous pouvons trouver dans la force irrésistible des exemples de vertu que nous offrent nos frères.

La Congrégation se voyant donc réduite par les circonstances à la triste nécessité de renoncer à conserver son préfet, ne crut pas que parmi ses membres, personne pût mieux le remplacer que son ami et son compagnon de souffrance M. Bertrand du Coin. C'était échanger un confesseur contre un autre confesseur ; mettre à sa tête un défenseur du St Siège, à la place d'un autre défenseur du Saint Siège ; c'était se cramponner à la chaire de St Pierre, se fixer invariablement sur ce fondement solide, comme sur un boulevard inexpugnable.

\section*{M. du Coin préfet}
% Mention marginale de la page manuscrite 292.
M. du Coin fut donc élu à une très forte majorité le jour de l'Assomption et prit possession de sa nouvelle dignité le jour de la Nativité de la Ste Vierge, 1814. Les circonstances ne permirent malheureusement pas qu'il pût gouverner pendant
% Source : PDF original, vue 146, côté droit, page manuscrite 293.
longtemps la Congrégation. Il profita néanmoins de la courte durée de son administration pour lui rendre de l'activité. Les sections étaient plus exactes que jamais à s'assembler ; elles avaient repris le cours de leurs bonnes œuvres ; les réceptions devenant de plus en plus fréquentes, bientôt la Congrégation se retrouva plus nombreuse qu'elle n'avait jamais été. Chaque nouvelle fête nous réunissait aux Chartreux ; les communions étaient nombreuses ; tout était parfaitement rentré dans l'état normal ; il ne manquait plus qu'une chose pour que notre organisation fût parfaitement rétablie. La Congrégation n'avait point de directeur Spirituel ; elle marchait toujours de son propre mouvement. Comme elle l'avait fait depuis le commencement de la persécution, ses autorités se bornaient dans les circonstances délicates, à aller prendre conseil auprès d'ecclésiastiques éclairés, que dans tous les temps il a été facile de rencontrer à Lyon.

On pourrait peut-être penser qu'une fois le calme revenu nous eussions dû aller nous remettre sous la main de M. Bochard que, dans le temps, l'Archevêque avait chargé de nous, et qui était toujours vicaire général. Mais d'après ce qui s'était passé entre lui et nous, cela n'eut guère été possible ; il est même douteux qu'il eut voulu y consentir, et d'un autre côté la Congrégation ne se fût pas du tout trouvé à l'aise avec lui. On pourrait avec vérité dire qu'entre lui et notre société, il y avait eu divorce pour cause d'incompatibilité d'humeur.

Ce qui pouvait nous rassurer dans cette position équivoque vis à vis de l'autorité Spirituelle, c'est que nous
% Source : PDF original, vue 147, côté gauche, page manuscrite 294.
étions toujours demeurés dans les meilleurs termes avec M. Courbon, premier vicaire général. C'était à lui que nous avions recours toutes les fois que pour solemniser nos fêtes, ou pour tout autre objet une permission de l'autorité nous était nécessaire, et il se prêtait toujours avec une grande bonté à seconder tous nos désirs. Toutefois il nous voyait avec peine livrés à nous-même et dénués de tout guide spirituel. Voyez, Messieurs, nous disait-il un jour dans toute la bonhomie, avec son ton un peu goguenard et son langage familier, Voyez, Messieurs, je suis obligé, à cause de vous, de me prendre tous les jours aux cheveux avec M. Bochard.

D'autre fois, nous parlant plus sérieusement il nous adressait ces paroles ou d'autres semblables : J'ai la plus grande confiance en vous, mais malgré votre zèle et toutes vos bonnes dispositions dont je ne doute pas et dont je ne pourrais pas douter puisque je les connais parfaitement, je n'en tremble pas moins sur ce qui pourrait arriver plus tard, si vous demeuriez dans l'état où vous êtes. Lorsque je pense que Pierre Valdo, qui a formé la secte des Vaudois, était de Lyon, je ne serai parfaitement tranquille à cet égard que lorsque nous aurons pu vous remettre entre les mains d'un ecclésiastique capable de vous diriger. Nous n'eussions certainement pas mieux demandé que de nous soumettre à l'exécution de son désir ; mais la seule difficulté était de trouver un prêtre capable, suffisamment indépendant et possédant d'ailleurs les qualités nécessaires pour remplir ce ministère auprès de nous. Or ce choix à faire était bien délicat ; en attendant qu'il eut pu y réussir, M. Courbon nous proposa
% Source : PDF original, vue 147, côté droit, page manuscrite 295.
d'accepter comme directeur, au moins pour la forme, M. l'abbé de Lupé. Vous le connaissez, dit-il, et vous savez bien qu'avec lui vous êtes toujours bien sûrs de mener vos affaires comme vous l'entendrez. Nous connaissions si bien ce bon et vénérable M. de Lupé, qu'il était lui-même Congréganiste. Il avait travaillé avec beaucoup de zèle pendant plusieurs années dans nos sections avant d'entrer dans le sacerdoce. Depuis qu'il était prêtre, il n'avait pas cessé d'être constamment et fidèlement attaché à la Congrégation. Nous fûmes donc heureux de voir ainsi toutes les difficultés s'aplanir, et nous acceptâmes cette offre avec reconnaissance ; mais l'excellent M. de Lupé dont le temps était alors entièrement consacré à d'autres fonctions, ne s'est jamais considéré comme directeur de la Congrégation ; il se contentait d'en porter le nom pour nous couvrir ainsi de son égide, et toutes les fois qu'il le pouvait, il se chargeait à nos jours de fête de nous donner la messe et de nous évangéliser. Quant aux assemblées générales, il ne pouvait y assister que fort rarement.

\section*{M. de Lupé directeur de la Congrégation}
% Mention marginale de la page manuscrite 294 ; placée après le récit de sa nomination.
La marche de la Congrégation paraissait de plus en plus ferme et assurée, lorsque le fatal débarquement de Napoléon vint replonger la France dans tous les troubles et les anxiétés inséparables d'une nouvelle révolution. Cet événement suspendit par contre coup toutes les relations des Congréganistes entre eux d'une manière plus absolue peut-être que pendant la longue période de la persécution dont nous venions à peine de sortir. L'action du gouvernement de la Congrégation fut si complètement arrêtée, et les Congréganistes se trouvèrent tellement dispersés par l'orage, qu'au troisième dimanche de
% Source : PDF original, vue 148, côté gauche, page manuscrite 296.
juillet 1815, il fut impossible de les réunir, même pour célébrer la fête par une simple communion générale. C'était la première fois depuis la fondation que ce grand jour passait ainsi inaperçu. Grâce à Dieu une circonstance aussi triste ne s'est pas renouvellée depuis.

\section*{Cent-jours}
% Mention marginale de la page manuscrite 296.
Le Seigneur eut pitié de nous. Marie jeta sur la France et peut-être aussi sur notre société, un regard de compassion. Encore une fois elle obtint notre grâce. Si la tempête avait été vive, heureusement elle fut courte. Des jours plus sereins reparurent ; les Congréganistes revinrent promptement reprendre leurs postes ; ils furent bientôt rassemblés ; mais l’orage leur avait enlevé leur Préfet. Ainsi que nous l'avons dit ailleurs, M. du Coin, qui s'était empressé de suivre le Roi à Gand en qualité de volontaire royal, s'était à la suite de cette expédition, déterminé à rentrer dans le service militaire. La Congrégation se voyait donc privée de sa présence ; il lui restait seulement pour l'édification des Congréganistes, le souvenir des vertus dont il nous avait rendu les témoins, et tous les bons exemples que nous offrirait encore sa conduite toujours également noble et chrétienne. Il restait aussi à tous les membres de la Congrégation l'espoir bien fondé de retirer de la communauté des biens spirituels établie entre nous, leur part dans les mérites que ce héros de la religion devait encore acquérir.

Le premier assistant continua à gouverner au nom du Préfet la Congrégation, jusqu'à l'époque assez rapprochée de l'élection de son successeur. Ce successeur fut M. Matthieu Garnier ; il fut élu le 15 août, jour de l'Assomption 1815 ; il fut installé dans la place où l'avait appelé le vœu de la
% Source : PDF original, vue 148, côté droit, page manuscrite 297, avant le chapitre quatorzième.
Congrégation, le 8 septembre, jour de la Nativité de la Ste Vierge de la même année.
